% =============================================================================
% Abstract
% One page maximum. Two or three paragraphs. Should answer:
%   1. What problem does the project address?
%   2. What approach was taken?
%   3. What is the headline result / contribution?
% Do NOT cite anything in the abstract.
% =============================================================================
\cleardoublepage
\chapter*{Abstract}
\thispagestyle{anonymous}
\addcontentsline{toc}{chapter}{Abstract}

Continual reinforcement learning seeks to train a single agent across a
sequence of related tasks without erasing earlier skills.
Regularisation based methods such as Elastic Weight Consolidation (EWC),
L2 anchoring and Memory Aware Synapses (MAS) are the standard
baselines on benchmarks such as the Doom-based COOM benchmark, but
their reported effectiveness has only been established under generous
training budgets that are unavailable on consumer hardware.

This dissertation investigates the inverse regime that arises in
resource constrained deployments: each task is allocated only one tenth
of the training steps reported in the original COOM benchmark. Under
this constraint the project compares EWC, L2 and MAS against an
unregularised Fine Tuning control on the CO4 task sequence of COOM, on
a consumer grade Apple Silicon CPU. EWC is run with three random
seeds; the remaining methods are single seed. Hyperparameters are
inherited from the COOM paper.

The headline finding is that all three regularisers
\emph{underperform} the Fine Tuning control by a factor of three to
seven on average training success rate, with the gap concentrated on
the final task in the sequence. Inspection of the regularisation
losses confirms that the Fisher and importance penalties activate as
designed; the failure is therefore a trade-off rather than an
implementation defect. The dissertation diagnoses the underlying
mechanism --- a ``Fisher on noise'' effect in which the source task
curvature estimate encodes the optimisation noise of an
under-converged source task rather than learned competence --- and
argues for compute aware continual learning methods that schedule
regularisation strength by the realised training budget.

As a direct test of this diagnosis the dissertation proposes
\emph{Online EWC}: a one-line modification of EWC that
replaces the additive Fisher accumulation across tasks with
an exponential moving average. The variant is run head to
head against the original EWC at two budgets. At the CPU
$20{,}000$-step budget, Online EWC at the conventional
decay $\gamma = 0.95$ underperforms the original EWC ($0.021$
vs $0.037$ on the CO4 average); at the paper-budget
$200{,}000$ steps run on a rented RTX 3090, the same
$\gamma = 0.95$ setting \emph{out}performs the original EWC
on the held-out CO4 average ($0.296$ vs $0.262$), with the
improvement concentrated on the task original EWC forgets
catastrophically. The reversal is exactly what the
Fisher-on-noise diagnosis predicts: at low budget the
Fisher captures noise that the EMA preserves; at paper
budget the Fisher captures the converged source-task
optimum that the EMA usefully decays. The diagnosis is
therefore supported by a falsifiable head-to-head, not
just by a qualitative argument. Multi-seed extension and
a denser $\gamma$ sweep are reported as the natural
follow-on work.

\medskip
\noindent\textbf{Keywords:} continual learning, reinforcement learning,
catastrophic forgetting, COOM, ViZDoom, Elastic Weight Consolidation,
plasticity--stability trade-off.
